\section{Scoring y desempates}

\subsection{Descripcion}

Este documento define el calculo oficial de puntajes del torneo, el uso de multiplicadores, la integracion de predicciones especiales y el orden de desempates del leaderboard.

\subsection{Alcance}

Aplica a:

\begin{itemize}
  \item puntaje por partido;
  \item puntaje acumulado del torneo;
  \item predicciones especiales de Campeon y Goleador;
  \item ordenamiento del leaderboard.
\end{itemize}

\subsection{Definiciones}

\begin{itemize}
  \item \textbf{Puntaje base}: suma de puntos obtenidos antes de aplicar multiplicador.
  \item \textbf{Puntaje acumulado}: suma total de puntajes validos del participante en el torneo.
  \item \textbf{Desempate}: regla de ordenamiento entre participantes con el mismo puntaje acumulado.
  \item \textbf{Desglose}: separacion visible del origen de puntos sin alterar el total oficial.
\end{itemize}

\subsection{Reglas}

\subsubsection{Puntaje por partido}

\begin{itemize}
  \item Acertar el outcome otorga 3 puntos.
  \item Acertar el marcador exacto otorga 2 puntos adicionales.
  \item El puntaje base maximo por partido es 5 puntos.
\end{itemize}

\subsubsection{Multiplicadores por fase}

\begin{center}
\begin{tabular}{lr}
\toprule
Fase & Multiplicador \\
\midrule
Grupos & x1.0 \\
Octavos de final & x1.25 \\
Cuartos de final & x1.5 \\
Semifinales & x1.75 \\
Tercer lugar & x1.75 \\
Final & x2.0 \\
\bottomrule
\end{tabular}
\end{center}

\begin{itemize}
  \item El multiplicador se aplica al puntaje del partido conforme a la fase oficial del torneo.
  \item El calculo conserva precision exacta al aplicar el multiplicador.
  \item La logica interna de scoring no debe redondear el valor calculado.
  \item La representacion visual del puntaje puede formatearse con precision controlada, pero ese formato no modifica el valor oficial de calculo.
\end{itemize}

\subsubsection{Predicciones especiales}

\begin{itemize}
  \item Campeon correcto otorga 20 puntos.
  \item Goleador correcto otorga 15 puntos.
  \item Las predicciones especiales forman parte del puntaje total acumulado del torneo.
  \item Las predicciones especiales no crean un leaderboard separado por defecto.
\end{itemize}

\subsubsection{Leaderboard}

\begin{itemize}
  \item El leaderboard se ordena por puntaje total acumulado.
  \item El puntaje total acumulado integra:
    \begin{itemize}
      \item puntaje de picks por partido;
      \item puntaje de predicciones especiales.
    \end{itemize}
  \item El sistema debe poder mostrar desglose del origen de puntos sin alterar el total oficial.
  \item Las vistas esperadas del leaderboard son:
    \begin{itemize}
      \item general;
      \item por jornada;
      \item por fase.
    \end{itemize}
  \item La vista general usa el total acumulado del torneo como valor principal.
  \item Las vistas filtradas por jornada o fase muestran como valor principal solo el puntaje atribuible al scope filtrado.
  \item Las predicciones especiales \texttt{Campeon} y \texttt{Goleador} no forman parte del puntaje principal de una vista filtrada por jornada o fase.
\end{itemize}

\subsubsection{Desempates}

Orden oficial de desempates:

\begin{enumerate}
  \item mayor cantidad de marcadores exactos acertados;
  \item mayor cantidad de outcomes acertados;
  \item mayor puntaje acumulado en fase KO;
  \item fecha y hora de registro mas temprana.
\end{enumerate}

\subsection{Edge Cases}

\begin{itemize}
  \item Si un resultado oficial se corrige, el puntaje afectado debe recalcularse de forma deterministica.
  \item Si un partido fue cancelado sin resultado oficial utilizable, ese partido no altera el ranking.
  \item Si el formato visual del puntaje usa redondeo o truncamiento para mostrar datos, ese formato no puede alterar el orden real del leaderboard.
  \item Si un puntaje es entero, la interfaz puede mostrarlo sin decimales.
  \item Si un puntaje tiene componente fraccionario, la interfaz debe mostrar exactamente dos decimales visibles.
\end{itemize}

\subsection{Invariantes}

\begin{itemize}
  \item Dos participantes con el mismo puntaje total deben ordenarse siempre con el mismo algoritmo de desempate.
  \item El recalculo no puede producir duplicacion de puntos.
  \item El leaderboard refleja solo resultados oficiales vigentes.
  \item El total acumulado es la unica base oficial de ranking.
  \item Las predicciones especiales suman al total, aunque puedan mostrarse en desglose.
  \item El formato visual del score nunca redefine ni altera el valor oficial de calculo.
\end{itemize}

\subsection{Presentacion visual acordada}

\begin{itemize}
  \item En leaderboard, el puntaje visible del scope activo es el dato primario.
  \item El desglose se presenta como contexto secundario, no como reemplazo del total oficial.
  \item En vista general, el desglose compacto puede separar \texttt{Partidos} y \texttt{Especiales}.
  \item El desglose expandido de rendimiento vive en perfil o detalle si el backend expone ese contrato.
\end{itemize}
